\section{Limitations}
\label{sec:limitations}

\begin{itemize}
  \item \textbf{Scale of our own training.} Our fine-tunes are 0.36B and 0.5B
  parameters with about two million tokens each. The measurement results
  hold on released models up to 8B, but we say nothing about how training
  itself behaves at that scale.
  \item \textbf{Behaviour on one model.} The evaluations, the LoRA sweep, DPO
  and the overlap analysis are all on Qwen2.5-0.5B and single-seed. SmolLM2
  replicates the spectral results only.
  \item \textbf{Released models are black boxes.} We do not know their
  post-training data or recipes, so they test the measurement but not any
  claim about alignment specifically.
  \item \textbf{Loss masking differs.} SFT computes loss on answer tokens only,
  the controls on every token. Tokens seen are matched; tokens trained on are
  not.
  \item \textbf{The cutoff test is narrow.} One model, one benchmark, one
  sampling mode. Gains reported by \citet{wu2025timber} on other benchmarks
  are untested here.
  \item \textbf{LoRA learning rates.} The lowest learning rate in the sweep
  won at every rank, so the rank-versus-GSM8K trend may partly reflect
  under-tuned learning rates.
\end{itemize}

\section{Conclusion}
\label{sec:conclusion}

Measures that count directions, whether by threshold, by energy share or by
entropy, describe fine-tuning updates as nearly full rank and put them within
a few percent of random matrices of the same shape. Stable rank shows the same
updates to be several times more concentrated than noise. Both descriptions
are accurate, because the updates combine a noise-like tail with a few
dominant directions, and the tail matters to the model. The practical lesson
is simple: a rank measure of a weight update means little without the same
measure on a matrix with no structure, and a claim about alignment needs a
fine-tune without alignment content as its control. With both in place, the
concentration of the update turns out to be a property of fine-tuning in
general, and an unchanged effective rank after post-training turns out to be
what any small update produces.

\subsubsection*{Broader Impact Statement}
This work is about measurement methodology for analysing fine-tuned language
models. It does not introduce new capabilities or training methods. Its main
practical consequence is caution about reading weight-space statistics as
evidence that alignment is shallow or deep; claims of either kind can affect
how much trust is placed in alignment training, and they should rest on
measurements with proper baselines.

% Reproducibility: code, configs, exact training rows, result CSVs and the
% saved singular-value spectra will be released with the camera-ready version
% (kept out of the anonymous submission).
